\chapter{ОЦЕНКА И АННОТАЦИЯ}
\label{app:evaluation}

\section{Состав серии}

\begin{table}[H]
    \centering
    \caption{Количество описаний и генераций}
    \label{tab:case-counts}
    \small
    \begin{tabular}{lccc}
        \toprule
        Набор & Описаний & Генераций на описание & Всего генераций \\
        \midrule
        acceptance2 & 8 & 3 & 24 \\
        acceptance3 & 8 & 3 & 24 \\
        acceptance4 & 6 & 3 & 18 \\
        composition & 4 & 1 & 4 \\
        \midrule
        Итого & 26 & --- & 70 \\
        \bottomrule
    \end{tabular}
\end{table}

\section{Behavior oracles}

Используются проверки сохранения суммы, установившегося значения, аналитической экспоненты, монотонности, числа смен знака, каскадного пика, задержанного ответа, изменения после события и допустимого диапазона. Для стационарного состояния дополнительная динамика не требуется; для остальных моделей хотя бы один stock должен меняться.

\begin{table}[H]
    \centering
    \caption{Уровни оценки}
    \label{tab:metric-levels}
    \small
    \begin{tabularx}{\textwidth}{L{3cm}X}
        \toprule
        Уровень & Содержание \\
        \midrule
        executable & Модель принята и симуляция выполнена \\
        finite & Все значения рядов конечны \\
        full horizon & Последняя точка совпадает с запрошенным final time \\
        semantic & Совпадает archetype, отсутствуют rule-based semantic errors, units/topology/graph пройдены \\
        behavior & Semantic, finite, full horizon, требуемая динамика и oracle выполнены \\
        \bottomrule
    \end{tabularx}
\end{table}

\section{Протокол экспертной аннотации}

Для будущей независимой оценки каждый эксперт должен ответить на следующие вопросы без доступа к названию варианта pipeline:

\begin{enumerate}
    \item Соответствуют ли выбранные stocks и flows описанию системы?
    \item Верны ли направления потоков и полярность связей?
    \item Согласованы ли единицы параметров и выражений?
    \item Соответствуют ли уравнения и начальные значения описанию?
    \item Допустимы ли сделанные допущения и указаны ли они?
    \item Совместимо ли поведение с ожидаемым качественным режимом?
    \item Пригодна ли модель как прототип для последующей экспертной доработки?
\end{enumerate}

Ответы фиксируются по заранее определённой шкале. Разногласия разрешает третий эксперт или согласованная adjudication. Автоматические semantic/behavior метрики не подменяются этой экспертной оценкой, а сопоставляются с ней.

\section{Blind-выгрузка и согласованность}

Для каждой метрики \codeid{semantic_score} и \codeid{behavior_score} эксперт выбирает \codeid{0}, \codeid{1} или \codeid{NA}. Последнее означает, что доступных материалов недостаточно для оценки; оно не засчитывается автоматически как успех или ошибка. Каждое задание содержит исходное описание, канонический ModelDocument и сводку симуляции, но не раскрывает ветвь генерации или служебные результаты проверки. Эксперту передаётся только каталог \codeid{expert/}; соответствия случайных идентификаторов исходным случаям сохраняются отдельно в \codeid{private/answer_key.json}.

Обе независимые разметки сначала используются для расчёта доли совпадений и Cohen's kappa, и только затем выполняется adjudication. По умолчанию задания без модели исключаются из расчёта согласованности: автоматический отказ генерации не должен искусственно увеличивать согласие экспертов. Такие задания можно включить отдельной опцией для анализа чувствительности.

На момент текущего эксперимента пакеты для независимой разметки сформированы, но два независимых эксперта и adjudication ещё не участвовали. Поэтому коэффициент согласованности не рассчитан, а приводимые в работе semantic/behavior результаты являются автоматическими проверками, а не экспертной оценкой.
